%% LyX 2.4.4 created this file.  For more info, see https://www.lyx.org/.
%% Do not edit unless you really know what you are doing.
\documentclass[english,handout]{beamer}
\usepackage[T1]{fontenc}
\usepackage[latin9]{inputenc}
\setcounter{secnumdepth}{3}
\setcounter{tocdepth}{3}
\usepackage{units}
\usepackage{graphicx}

\makeatletter
%%%%%%%%%%%%%%%%%%%%%%%%%%%%%% Textclass specific LaTeX commands.
% this default might be overridden by plain title style
\newcommand\makebeamertitle{\frame{\maketitle}}%
% (ERT) argument for the TOC
\AtBeginDocument{%
  \let\origtableofcontents=\tableofcontents
  \def\tableofcontents{\@ifnextchar[{\origtableofcontents}{\gobbletableofcontents}}
  \def\gobbletableofcontents#1{\origtableofcontents}
}

%%%%%%%%%%%%%%%%%%%%%%%%%%%%%% User specified LaTeX commands.
\beamertemplateshadingbackground{red!5}{structure!5}

\usepackage{beamerthemeshadow}
\usepackage{pgfnodes,pgfarrows,pgfheaps}
\usepackage{caption}

\beamertemplatetransparentcovereddynamicmedium

\pgfdeclareimage[width=0.6cm]{icsi-logo}{beamer-icsi-logo}
\logo{\pgfuseimage{icsi-logo}}

\newcommand{\Class}[1]{\operatorname{\mathchoice
  {\text{\small #1}}
  {\text{\small #1}}
  {\text{#1}}
  {\text{#1}}}}

\newcommand{\Lang}[1]{\operatorname{\text{\textsc{#1}}}}

% This gets defined by beamerbasecolor.sty, but only at the beginning of
% the document
\colorlet{averagebackgroundcolor}{normal text.bg}

\newcommand{\tape}[3]{%
  \color{structure!30!averagebackgroundcolor}
  \pgfmoveto{\pgfxy(-0.5,0)}
  \pgflineto{\pgfxy(-0.6,0.1)}
  \pgflineto{\pgfxy(-0.4,0.2)}
  \pgflineto{\pgfxy(-0.6,0.3)}
  \pgflineto{\pgfxy(-0.4,0.4)}
  \pgflineto{\pgfxy(-0.5,0.5)}
  \pgflineto{\pgfxy(4,0.5)}
  \pgflineto{\pgfxy(4.1,0.4)}
  \pgflineto{\pgfxy(3.9,0.3)}
  \pgflineto{\pgfxy(4.1,0.2)}
  \pgflineto{\pgfxy(3.9,0.1)}
  \pgflineto{\pgfxy(4,0)}
  \pgfclosepath
  \pgffill

  \color{structure}  
  \pgfputat{\pgfxy(0,0.7)}{\pgfbox[left,base]{#1}}
  \pgfputat{\pgfxy(0,-0.1)}{\pgfbox[left,top]{#2}}

  \color{black}
  \pgfputat{\pgfxy(-.1,0.25)}{\pgfbox[left,center]{\texttt{#3}}}%
}

\newcommand{\shorttape}[3]{%
  \color{structure!30!averagebackgroundcolor}
  \pgfmoveto{\pgfxy(-0.5,0)}
  \pgflineto{\pgfxy(-0.6,0.1)}
  \pgflineto{\pgfxy(-0.4,0.2)}
  \pgflineto{\pgfxy(-0.6,0.3)}
  \pgflineto{\pgfxy(-0.4,0.4)}
  \pgflineto{\pgfxy(-0.5,0.5)}
  \pgflineto{\pgfxy(1,0.5)}
  \pgflineto{\pgfxy(1.1,0.4)}
  \pgflineto{\pgfxy(0.9,0.3)}
  \pgflineto{\pgfxy(1.1,0.2)}
  \pgflineto{\pgfxy(0.9,0.1)}
  \pgflineto{\pgfxy(1,0)}
  \pgfclosepath
  \pgffill

  \color{structure}  
  \pgfputat{\pgfxy(0.25,0.7)}{\pgfbox[center,base]{#1}}
  \pgfputat{\pgfxy(0.25,-0.1)}{\pgfbox[center,top]{#2}}

  \color{black}
  \pgfputat{\pgfxy(-.1,0.25)}{\pgfbox[left,center]{\texttt{#3}}}%
}

\pgfdeclareverticalshading{heap1}{\the\paperwidth}%
  {color(0pt)=(black); color(1cm)=(structure!65!white)}
\pgfdeclareverticalshading{heap2}{\the\paperwidth}%
  {color(0pt)=(black); color(1cm)=(structure!55!white)}
\pgfdeclareverticalshading{heap3}{\the\paperwidth}%
  {color(0pt)=(black); color(1cm)=(structure!45!white)}
\pgfdeclareverticalshading{heap4}{\the\paperwidth}%
  {color(0pt)=(black); color(1cm)=(structure!35!white)}
\pgfdeclareverticalshading{heap5}{\the\paperwidth}%
  {color(0pt)=(black); color(1cm)=(structure!25!white)}
\pgfdeclareverticalshading{heap6}{\the\paperwidth}%
  {color(0pt)=(black); color(1cm)=(red!35!white)}

\newcommand{\heap}[5]{%
  \begin{pgfscope}
    \color{#4}
    \pgfheappath{\pgfxy(0,#1)}{\pgfxy(-#2,0)}{\pgfxy(#2,0)}
    \pgfclip
    \begin{pgfmagnify}{1}{#1}
      \pgfputat{\pgfpoint{-.5\paperwidth}{0pt}}{\pgfbox[left,base]{\pgfuseshading{heap#5}}}
    \end{pgfmagnify}
  \end{pgfscope}
  %\pgffill
  
  \color{#4}
  \pgfheappath{\pgfxy(0,#1)}{\pgfxy(-#2,0)}{\pgfxy(#2,0)}
  \pgfstroke

  \color{white}
  \pgfheaplabel{\pgfxy(0,#1)}{#3}%
}


\pgfdeclareradialshading{graphnode}
  {\pgfpoint{-3pt}{3.6pt}}%
  {color(0cm)=(beamerexample!15);
    color(2.63pt)=(beamerexample!75);
    color(5.26pt)=(beamerexample!70!black);
    color(7.6pt)=(beamerexample!50!black);
    color(8pt)=(beamerexample!10!averagebackgroundcolor)}

\newcommand{\graphnode}[2]{
  \pgfnodecircle{#1}[virtual]{#2}{8pt}
  \pgfputat{#2}{\pgfbox[center,center]{\pgfuseshading{graphnode}}}
}

%\setbeamerfont{author}{size=\Huge} 
%\setbeamerfont{date}{size=\small} 
%\setbeamerfont{exampletext}{series=\bfseries}
%\setbeamerfont{frametitle}{series=\bfseries,size=\large}
%\setbeamerfont{framesubtitle}{series=\mdseries,size=\large}
%\setbeamerfont{section}{series=\bfseries} 
\setbeamerfont{title}{series=\bfseries}

\setbeamercovered{invisible}

\usepackage{pgfpages}
\pgfpagesuselayout{4 on 1}[a4paper, landscape, border shrink=7mm]
\pgfpageslogicalpageoptions{1}{border code=\pgfusepath{stroke}}
\pgfpageslogicalpageoptions{2}{border code=\pgfusepath{stroke}}
\pgfpageslogicalpageoptions{3}{border code=\pgfusepath{stroke}}
\pgfpageslogicalpageoptions{4}{border code=\pgfusepath{stroke}}

\makeatother

\usepackage{babel}
\begin{document}
\title{Calculus II}
\subtitle{MTH 231~~|~~Section C~~|~~Fall 2012}
\institute{{\Large\textbf{Johnson C. Smith }}}
\author{Dr.\,James Ashe}
\date{{}}

\makebeamertitle
%%suppress auto date
\title{Calc II: 4.8 Antiderivaties}

\author{}

\institute{}

\subtitle{}

\makebeamertitle
\begin{frame}{Review}

\begin{exampleblock}{Using limits to avoid division by zero}

A ball rolls down a ramp so that its height $h$ from the top after
$t$ seconds is $t^{2}$. What is the instantaneous rate of change
of the $h$ with respect to $t$ after $3$ seconds?

\begin{itemize}
\item <2->i.e., what is its velocity at $t=3$.
\end{itemize}
\end{exampleblock}
\begin{columns}


\column{6.5cm}

\vspace{-.75cm}

\begin{align*}
\uncover<3->{\frac{\triangle h}{\triangle t}} & \uncover<3->{=\frac{(4)^{2}-3^{2}}{4-3}=7\mbox{ ft/sec}}\\
\uncover<4->{\frac{\triangle h}{\triangle t}} & \uncover<4->{=\frac{(3.25)-3^{2}}{3.25-3}=6.25\mbox{ ft/sec}}\\
\uncover<5->{\vdots}\\
\uncover<6->{{\displaystyle \lim_{t\rightarrow3}}\frac{\triangle h}{\triangle t}} & \uncover<6->{={\displaystyle \lim_{t\rightarrow3}\frac{t^{2}-3^{2}}{t-3}}={\displaystyle \lim_{t\rightarrow3}\frac{(t-3)(t+3)}{(t-3)}}}\\
\uncover<7->{\displaystyle \lim_{t\rightarrow3}\frac{\triangle h}{\triangle t}} & \uncover<7->{={\displaystyle \lim_{t\rightarrow3}}t+3=6}
\end{align*}


\column{5cm}

\includegraphics[width=5cm]{images/4-8_avg_ROC}<3|handout:0>\includegraphics[width=5cm]{images/4-8_avg_ROC_2}<4-6|handout:0>\includegraphics[width=5cm]{images/4-8_avg_ROC_3}<7>
\end{columns}
\end{frame}

\begin{frame}{}

\begin{columns}


\column{6.0cm}
\begin{block}{Notation: The derivative of a function, $f(x)$}

\begin{itemize}
\item <2->The derivative of $f$ with respect to $x$ written: $f'(x)=\frac{d}{dx}f(x)=\frac{df(x)}{dx}=\frac{df}{dx}$.
\item <5->The derivative of $f$ with respect to $x$ \emph{at a value
``$a$}'': $f'(a)=\frac{d}{dx}f(a)$.

\begin{itemize}
\item <6->This is the ``slope'' of $f$ at $a$.
\end{itemize}
\end{itemize}
\end{block}
\begin{exampleblock}<3->{Example}

\begin{itemize}
\item <3->$f(x)=2x^{3}-4x^{2}+4\Rightarrow$
\item <4->[]$f'(x)=6x^{2}-8x$
\item <7->For $a=$1, $f'(1)=6(1)^{2}-8(1)=-2$
\end{itemize}
\end{exampleblock}

\column{5.0cm}

\includegraphics[width=5cm]{images/function_2}<3|handout:0>\includegraphics[width=5cm]{images/function_1}<4-6|handout:0>\includegraphics[width=5cm]{images/function_3}<7->
\end{columns}
\end{frame}

\begin{frame}{}

\begin{exampleblock}{Questions}

\begin{enumerate}
\item <1->[1)]Suppose you are a driving at a constant rate of $50$ $\nicefrac{mi}{hr}$.
How far do you travel in the first $3$ hours? 

\begin{enumerate}
\item []\uncover<2->{$50\times3=150$} 
\end{enumerate}
\item <3->[2)]Suppose a ball rolls down a ramp at velocity of $2t$ $\nicefrac{ft}{sec}$
where $t$ is the number of seconds after it started to roll. How
far does the ball travel during the first $3$ seconds?

\begin{enumerate}
\item <4->[]Here the velocity varies according to time. Clearly, the derivative
is related to this problem, but instantaneous velocity occurs instantly.
This problem will again require limits, but instead of tangent lines
we will find areas.
\end{enumerate}
\end{enumerate}
\end{exampleblock}
\end{frame}

\begin{frame}{Antiderivatives}

\begin{definition}

A function $F$ is an \textbf{\emph{antiderivative }}of $f$ on an
interval $I$ if and only if $F'(x)=f(x)$ for all $x\in I$.

Find all the antiderivatives of the following functions
\end{definition}

\begin{enumerate}
\item <2->[ex 10)]$g(x)=11x^{10}$ \uncover<3->{$\Rightarrow$$G(x)=x^{11}+C$
where $C$ is some constant number.}
\end{enumerate}
\uncover<4->{Note that: $\frac{d}{dx}(x^{11}+2)=\frac{d}{dx}(x^{11}+\pi)=\frac{d}{dx}(x^{11}+\frac{1}{2})\dots=10x^{11}$.}

\uncover<5->{So $x^{11}+C$ represents the \emph{family} of antiderivatives
of $f$.}
\end{frame}

\begin{frame}{}

\begin{block}{Notation}

$\int f(x)dx$, the\emph{ indefinite integral}, represents the family
of antiderivatives of $f(x)$.
\end{block}
\uncover<2->{If $F(x)=\frac{x^{p+1}}{p+1}+C$ } \uncover<3->{then
$F'(x)=(p+1)\frac{1}{p+1}x^{p+1-1}=x^{p}$.} 

\uncover<4->{So then $F'(x)=x^{p}=f(x)$.}
\begin{block}<5->{Theorems: }

Power rule:$\int x^{p}dx=\frac{x^{P+1}}{p+1}+C$ where $p\neq-1$
is a real number.

Sum rule:$\int(f(x)+g(x))dx=\int f(x)dx+\int g(x)dx$

Constant rule: $\int cf(x)dx=c\int f(x)dx$
\end{block}
\uncover<6->{ex) $\int(-6\frac{1}{z^{7}}+2)dz=$}\uncover<7->{$-6\int z^{-7}dz+\int2dz$}

\hspace{1.5cm}\uncover<8->{$=z^{-6}+2z$}\uncover<9->{$+C$}
\end{frame}

\begin{frame}{}

Recall the chain rule, $(f\circ g)'(x)=f'(g(x))g'(x).$
\begin{itemize}
\item <2->So then $\frac{d}{dx}(\sin ax)=$\uncover<3->{\textrm{$a\cdot\cos(ax)$}}
\item <4->Similarly we can find the indefinite integrals of trigonometric
functions; see Table 4.5 pg. 243.
\end{itemize}
\end{frame}

\begin{frame}{}

\begin{exampleblock}{Introduction to Differential Equations}

A ball rolls down a ramp so that its velocity is $v(t)=2t$ feet per
second. The distance of the ball from the top of the ramp is 3 feet
after 1 second. Find the function giving the distance of the ball
from the top of the ramp at $t$ seconds. 
\end{exampleblock}
\begin{itemize}
\item []<2->$\int2tdt=t^{2}+C$; an infinite class of solutions.
\item []<3->$h(t)=t^{2}+C\Rightarrow h(1)=1+C=3\Rightarrow C=2$
\item []<4->So $h(t)=t^{2}+2$. 
\end{itemize}
\end{frame}

\begin{frame}{}

\textbf{HW: }15, 19, 24, 30, 38, 58

\vspace{1cm}

Next time, 5.1: Approximating areas under curves
\end{frame}

\end{document}
